%% kapa: exam=2 round=23 year=2012 period=1 subject=appraisal_practice subject_ko=감정평가실무 seq=3 number=3 points=20
%% source: Q-Net 감정평가사 기출문제 (cst00309) artlSeq=5056012 file=raw/1008/5056012/1교시(감정평가실무).hwp
%% conversion: kordoc 4.24.51 -> tools/gen2.py (auto); review=proofread-sonnet; flags=none
\begin{longquestion}{【문제 3】}
甲은 본인이 소유하고 있는 토지를 이용하여 공장을 신축하였다. B시청에서는 당해 사업이 개발부담금 부과대상사업에 해당되어 개발부담금을 부과하려고 한다. 주어진 자료를 활용하여 다음 물음에 답하시오. (20점) \par
\medskip\noindent 물음 1) 30-2번지에 대하여 개시시점지가와 종료시점지가를 산정하시오. (10점) \par
\medskip\noindent 물음 2) 30-4번지에 대하여 개시시점지가와 종료시점지가를 산정하시오. (5점) \par
\medskip\noindent 물음 3) 30-5번지에 대하여 개시시점지가(매입가액기준)와 종료시점지가를 산정하시오. (5점) \par
\textbf{<자료 1> 기본적 사항} \par
1. 개발사업 인가일: 2011.10.01 \par
2. 개발사업 준공인가일: 2012.08.30 \par
3. 사업인가조건: 30-2번지 중 일부(500m\textsuperscript{2})를 도로 등으로 기부채납 \par
4. 현장조사 완료일: 2012.09.09 \par
\textbf{<자료 2> 대상토지자료} \par
1. 기본내용 \par
{\small\centering
\begin{tabular}{|c|c|c|c|c|c|}
\hline
\makecell{일련번호} & \makecell{토지소재} & \makecell{지번} & \makecell{지목} & \makecell{이용상황} & \makecell{면적(m\textsuperscript{2})} \\
\hline
\makecell{①} & \makecell{B시 D동} & \makecell{30-2} & \makecell{전} & \makecell{전} & \makecell{3,500} \\
\hline
\makecell{②} & \makecell{B시 D동} & \makecell{30-4} & \makecell{답} & \makecell{답} & \makecell{3,000} \\
\hline
\makecell{③} & \makecell{B시 D동} & \makecell{30-5} & \makecell{답} & \makecell{답} & \makecell{1,000} \\
\hline
\end{tabular}\par}

2. 용도지역: 계획관리지역임 \par
3. 토지 특성: 일련번호 ①, ②, ③ 토지 특성은 동일함. \par
1) 개발 전: 세로가, 부정형, 완경사 \par
2) 개발 후: 소로한면, 세장형, 평지 \par
\textbf{<자료 3> 가격결정을 위한 참고자료} \par
1. 표준지공시지가 현황 \par
{\small\centering
\begin{tabular}{|c|c|c|c|c|c|c|c|c|c|}
\hline
\multirow{2}{*}{\makecell{기호}} & \multirow{2}{*}{\makecell{소재지}} & \multirow{2}{*}{\makecell{면적\newline{}(m\textsuperscript{2})}} & \multirow{2}{*}{\makecell{지목}} & \multirow{2}{*}{\makecell{이용\newline{}상황}} & \multirow{2}{*}{\makecell{용도지역}} & \multirow{2}{*}{\makecell{도로교통}} & \multirow{2}{*}{\makecell{형상\newline{}지세}} & \multicolumn{2}{|c|}{\makecell{공시지가(원/m\textsuperscript{2})}} \\
\cline{9-10}
 &  &  &  &  &  &  &  & \makecell{2011년} & \makecell{2012년} \\
\hline
\makecell{①} & \makecell{D동\newline{}32} & \makecell{500} & \makecell{답} & \makecell{답} & \makecell{계획관리\newline{}지역} & \makecell{세로\newline{}(가)} & \makecell{세장형 평지} & \makecell{50,000} & \makecell{55,000} \\
\hline
\makecell{②} & \makecell{D동 50-1} & \makecell{1,000} & \makecell{장} & \makecell{공업용} & \makecell{계획관리\newline{}지역} & \makecell{세로\newline{}(가)} & \makecell{세장형\newline{}평지} & \makecell{200,000} & \makecell{210,000} \\
\hline
\end{tabular}\par}

2. 개별공시지가 \par
{\small\centering
\begin{tabular}{|c|c|c|c|c|}
\hline
\makecell{일련번호} & \makecell{토지소재} & \makecell{지번} & \makecell{2011년(원/m\textsuperscript{2})} & \makecell{2012년(원/m\textsuperscript{2})} \\
\hline
\makecell{1} & \makecell{B시 D동} & \makecell{30-2} & \makecell{-} & \makecell{-} \\
\hline
\makecell{2} & \makecell{B시 D동} & \makecell{30-4} & \makecell{45,000} & \makecell{50,000} \\
\hline
\makecell{3} & \makecell{B시 D동} & \makecell{30-5} & \makecell{45,000} & \makecell{50,000} \\
\hline
\end{tabular}\par}

3. 甲은 30-5번지를 경매로 60,000원/m\textsuperscript{2}에 낙찰 받아 2011.06.10에 소유권 이전을 완료하였다. \par
4. 지가변동률(\%) \par
{\small\centering
\begin{tabular}{|c|c|c|c|}
\hline
\makecell{구분} & \makecell{A도 평균} & \makecell{B시 평균} & \makecell{B시 계획\newline{}관리지역} \\
\hline
\makecell{2011.01.01 ～ 2011.06.10} & \makecell{3.1} & \makecell{7.1} & \makecell{5.1} \\
\hline
\makecell{2011.01.01 ～ 2011.10.01} & \makecell{5.5} & \makecell{10.0} & \makecell{7.5} \\
\hline
\makecell{2011.06.10 ～ 2011.10.01} & \makecell{1.0} & \makecell{2.5} & \makecell{1.5} \\
\hline
\makecell{2011.01.01 ～ 2012.08.30} & \makecell{12.5} & \makecell{13.5} & \makecell{12.8} \\
\hline
\makecell{2011.10.01 ～ 2012.08.30} & \makecell{10.0} & \makecell{11.5} & \makecell{11.0} \\
\hline
\makecell{2012.01.01 ～ 2012.08.30} & \makecell{8.5} & \makecell{10.5} & \makecell{9.5} \\
\hline
\makecell{2011.01.01 ～ 2012.09.09} & \makecell{13.5} & \makecell{14.8} & \makecell{13.8} \\
\hline
\makecell{2012.01.01 ～ 2012.09.09} & \makecell{9.8} & \makecell{11.8} & \makecell{10.8} \\
\hline
\makecell{2011.10.01 ～ 2012.09.09} & \makecell{10.8} & \makecell{12.0} & \makecell{11.8} \\
\hline
\end{tabular}\par}

5. 지역요인: 동일함 \par
6. 개별요인 비교치(토지가격비준표와 동일) \par
1) 도로접면 \par
{\small\centering
\begin{tabular}{|c|c|c|}
\hline
\makecell{구 분} & \makecell{소로한면} & \makecell{세로가} \\
\hline
\makecell{소로한면} & \makecell{1.00} & \makecell{0.93} \\
\hline
\makecell{세로가} & \makecell{1.07} & \makecell{1.00} \\
\hline
\end{tabular}\par}

2) 형상 \par
{\small\centering
\begin{tabular}{|c|c|c|}
\hline
\makecell{구 분} & \makecell{세장형} & \makecell{부정형} \\
\hline
\makecell{세장형} & \makecell{1.00} & \makecell{0.96} \\
\hline
\makecell{부정형} & \makecell{1.04} & \makecell{1.00} \\
\hline
\end{tabular}\par}

3) 지세 \par
{\small\centering
\begin{tabular}{|c|c|c|}
\hline
\makecell{구 분} & \makecell{평지} & \makecell{완경사} \\
\hline
\makecell{평지} & \makecell{1.00} & \makecell{0.97} \\
\hline
\makecell{완경사} & \makecell{1.03} & \makecell{1.00} \\
\hline
\end{tabular}\par}

4) 이용상황 \par
{\small\centering
\begin{tabular}{|c|c|c|c|}
\hline
\makecell{구 분} & \makecell{전} & \makecell{답} & \makecell{공업용} \\
\hline
\makecell{전} & \makecell{1.00} & \makecell{0.97} & \makecell{1.33} \\
\hline
\makecell{답} & \makecell{1.03} & \makecell{1.00} & \makecell{1.39} \\
\hline
\makecell{공업용} & \makecell{0.75} & \makecell{0.72} & \makecell{1.00} \\
\hline
\end{tabular}\par}

7. 기타요인 산정을 위한 자료 \par
1) 평가선례 \par
{\small\centering
\begin{tabular}{|c|c|c|c|c|c|c|c|}
\hline
\makecell{기호} & \makecell{소재지} & \makecell{지목} & \makecell{면적\newline{}(m\textsuperscript{2})} & \makecell{이용상황} & \makecell{용도\newline{}지역} & \makecell{단가(원/m\textsuperscript{2})} & \makecell{가격시점} \\
\hline
\makecell{①} & \makecell{D동\newline{}33} & \makecell{답} & \makecell{300} & \makecell{답} & \makecell{계획관리\newline{}지역} & \makecell{65,000} & \makecell{2011.01.01} \\
\hline
\makecell{②} & \makecell{D동\newline{}51} & \makecell{장} & \makecell{500} & \makecell{공업용} & \makecell{계획관리\newline{}지역} & \makecell{240,000} & \makecell{2011.01.01} \\
\hline
\makecell{③} & \makecell{D동\newline{}35} & \makecell{전} & \makecell{500} & \makecell{전} & \makecell{계획관리\newline{}지역} & \makecell{75,000} & \makecell{2012.01.01} \\
\hline
\makecell{④} & \makecell{D동\newline{}52} & \makecell{장} & \makecell{1,000} & \makecell{공업용} & \makecell{계획관리\newline{}지역} & \makecell{270,000} & \makecell{2012.01.01} \\
\hline
\end{tabular}\par}

2) 지가변동률은 상기에서 제시한 자료와 동일하고 대상토지와 평가선례와의 지역요인 및 개별요인은 동일함. \par
\end{longquestion}
